\chapter{Solvable vs Perpetual Problems: Masalah yang Bisa vs Tidak Bisa Diselesaikan}

\begin{insightbox}
\textit{``69\% konflik pernikahan adalah perpetual --- tidak akan pernah terselesaikan sepenuhnya. Pasangan bahagia bukan yang tanpa konflik, melainkan yang tahu cara \highlight{mengelola} konflik dengan bijaksana. Rahasianya bukan memenangkan pertarungan, tapi memahami perangnya.''}
\end{insightbox}

\section{Paradigma Baru: Dua Dunia Masalah yang Berbeda}

Dr. John Gottman, setelah meneliti ribuan pasangan selama lebih dari 40 tahun, menemukan fakta yang mengejutkan: \textbf{kebanyakan pasangan salah memahami konflik mereka}. Mereka memperlakukan setiap perbedaan sebagai masalah yang harus ``dipecahkan'' --- padahal ada dua kategori masalah yang sangat berbeda dalam hubungan.

\begin{praktikbox}[Revolusi Pemikiran Gottman]
Sebelum Gottman, terapi pernikahan berfokus pada ``problem-solving'' --- mencari solusi untuk setiap konflik. Gottman membuktikan bahwa pendekatan ini:\
\begin{itemize}
    \item \textbf{Gagal} pada 69\% kasus karena masalahnya memang \textbf{tidak ada solusinya}
    \item Membuat pasangan frustrasi dan merasa gagal
    \item Mengabaikan makna simbolik di balik perbedaan
\end{itemize}
\textit{Solusinya: Bedakan dulu masalahnya, baru tentukan strateginya.}
\end{praktikbox}

\subsection{Mengapa Pemahaman Ini Menyelamatkan Hubungan}

Tanpa pemahaman ini:
\begin{itemize}
    \item Kamu menghabiskan energi mencari ``solusi'' yang tidak ada
    \item Merasa frustrasi karena pasangan ``tidak mengerti''
    \item Menyimpulkan bahwa hubunganmu ``gagal'' karena konflik terus berulang
    \item Akhirnya menyerah: ``Kita tidak cocok''
\end{itemize}

Dengan pemahaman ini:
\begin{itemize}
    \item Kamu \textbf{menerima} bahwa beberapa perbedaan adalah bagian dari hubungan
    \item Fokus pada \textbf{mengelola} konflik, bukan menghilangkannya
    \item Melihat perbedaan sebagai sumber \textbf{pertumbuhan}, bukan ancaman
    \item Membangun kedekatan dari pemahaman, bukan kesamaan
\end{itemize}

\section{69\% Statistic: Kenapa Kebanyakan Konflik Abadi}

\subsection{Penelitian Gottman yang Mengubah Segalanya}

Dalam studi longitudinal yang melacak pasangan selama 14 tahun, Gottman menemukan:

\begin{center}
\begin{tikzpicture}[
    stat/.style={draw=#1, line width=3pt, rounded corners=15pt,
                 fill=#1!10, minimum width=5cm, minimum height=3cm, align=center}
]
    \node[stat=accent] (perpetual) at (-4,0) {
        \textbf{\Huge 69\%}\\[0.3cm]
        \textbf{PERPETUAL}\\[0.2cm]
        Konflik yang tidak\\pernah terselesaikan
    };
    
    \node[stat=sagegreen] (solvable) at (4,0) {
        \textbf{\Huge 31\%}\\[0.3cm]
        \textbf{SOLVABLE}\\[0.2cm]
        Konflik yang bisa\\diselesaikan sepenuhnya
    };
    
    \draw[<->, line width=2pt, primary] (perpetual) -- (solvable) node[midway, above, font=\small] {Semua Konflik Pernikahan};
\end{tikzpicture}
\end{center}

\subsection{Mengapa 69\%? Dasar Psikologisnya}

Angka ini bukan kebetulan. Ini mencerminkan \textbf{fakta fundamental} tentang manusia:

\begin{enumerate}
    \item \textbf{Kepribadian Stabil:} Trait dasar seperti ekstroversi, neuroticism, openness stabil sepanjang hidup
    \item \textbf{Nilai Terbentuk di Masa Kecil:} Konsep tentang uang, waktu, kebersihan, parenting tertanam dalam
    \item \textbf{Attachment Style:} Cara kita terhubung dengan orang lain sulit diubah total
    \item \textbf{Kebutuhan Dasar Berbeda:} Setiap orang punya level kebutuhan akan kedekatan, ruang, struktur yang berbeda
\end{enumerate}

\begin{tipbox}[Implikasi Praktis 69\%]
\begin{itemize}
    \item Jika kalian berkonflik tentang \textbf{topik yang sama} lebih dari 3 kali, kemungkinan besar itu \textbf{perpetual problem}
    \item Berhenti mencari ``solusi akhir'' --- itu seperti mencari kota fata morgana
    \item Alihkan fokus: dari ``bagaimana menyelesaikan'' ke ``bagaimana mengelola dengan baik''
\end{itemize}
\end{tipbox}

\subsection{Perbedaan Kunci dalam Perspektif}

\begin{insightbox}
\textbf{Paradoks Konflik:} Pasangan yang bahagia dan pasangan yang tidak bahagia memiliki \textbf{jumlah perpetual problem yang sama}. Perbedaannya bukan pada \textit{ada tidaknya} konflik, tapi pada \textit{cara menjalani} konflik tersebut.
\end{insightbox}

\section{Tabel Karakteristik Komprehensif: Solvable vs Perpetual}

Berikut adalah perbandingan detail lebih dari 10 dimensi kunci:

\begin{center}
\renewcommand{\arraystretch}{1.6}
\begin{tabular}{|>{\raggedright}p{3.2cm}|>{\raggedright}p{5.8cm}|>{\raggedright\arraybackslash}p{5.8cm}|}
\hline
\rowcolor{primary!20}
\textbf{Dimensi} & \textbf{SOLVABLE PROBLEMS} & \textbf{PERPETUAL PROBLEMS} \\
\hline
\textbf{1. Sifat Masalah} & Situasional dan konkret & Fundamental dan karakterologis \\
\hline
\textbf{2. Frekuensi} & Muncul sekali atau jarang & Berulang terus menerus \\
\hline
\textbf{3. Konteks} & Spesifik pada situasi tertentu & Muncul di berbagai konteks \\
\hline
\textbf{4. Makna Simbolik} & Minimal atau tidak ada & Sangat kaya dan dalam \\
\hline
\textbf{5. Solusi} & Ada solusi konkret dan clear & Tidak ada solusi ``benar'' \\
\hline
\textbf{6. Emosi Terlibat} & Emosi moderate dan terkontrol & Emosi intens dan personal \\
\hline
\textbf{7. Sejarah} & Terbatas pada hubungan ini & Sering terbawa dari masa lalu \\
\hline
\textbf{8. Area Terkait} & Satu area spesifik & Menyangkut banyak area hidup \\
\hline
\textbf{9. Respon Pasangan} & Kooperatif dalam problem-solving & Defensif dan protektif \\
\hline
\textbf{10. Potensi Resolusi} & 90-100\% bisa diselesaikan & 0\% diselesaikan, 100\% di-manage \\
\hline
\textbf{11. Efek Jangka Panjang} & Tidak mempengaruhi fondasi hubungan & Membentuk ``kebohongan jarak'' \\
\hline
\textbf{12. Kebutuhan Skill} & Problem-solving komunikasi & Dialogue dan empathy mendalam \\
\hline
\textbf{13. Tanda Sukses} & Masalah hilang/pergi & Masalah ada tapi tidak merusak \\
\hline
\textbf{14. Ancaman Hubungan} & Rendah jika diselesaikan & Tinggi jika menjadi gridlock \\
\hline
\end{tabular}
\end{center}

\subsection{Visual Perbandingan}

\begin{center}
\begin{tikzpicture}[
    solvable/.style={draw=sagegreen, line width=2pt, rounded corners=10pt, fill=sagegreen!10, minimum width=6cm, minimum height=5cm},
    perpetual/.style={draw=accent, line width=2pt, rounded corners=10pt, fill=accent!10, minimum width=6cm, minimum height=5cm}
]
    % Solvable Box
    \node[solvable] (sol) at (-4,0) {};
    \node[font=\Large\bfseries, sagegreen] at (-4,2) {SOLVABLE};
    \node[align=left, text width=5.5cm, font=\small] at (-4,0) {
        \faCheck\ Masalah spesifik\\[0.2cm]
        \faCheck\ Ada solusi konkret\\[0.2cm]
        \faCheck\ Emosi moderate\\[0.2cm]
        \faCheck\ Situasional\\[0.2cm]
        \faCheck\ Dapat diselesaikan
    };
    
    % Perpetual Box
    \node[perpetual] (per) at (4,0) {};
    \node[font=\Large\bfseries, accent] at (4,2) {PERPETUAL};
    \node[align=left, text width=5.5cm, font=\small] at (4,0) {
        \faExclamationTriangle\ Masalah berulang\\[0.2cm]
        \faExclamationTriangle\ Tidak ada solusi\\[0.2cm]
        \faExclamationTriangle\ Emosi intens\\[0.2cm]
        \faExclamationTriangle\ Fundamental\\[0.2cm]
        \faExclamationTriangle\ Hanya bisa di-manage
    };
\end{tikzpicture}
\end{center}

\section{Contoh Perbandingan dalam Kehidupan Nyata}

\begin{center}
\renewcommand{\arraystretch}{1.5}
\begin{tabular}{|>{\raggedright}p{2.5cm}|>{\raggedright}p{5cm}|>{\raggedright\arraybackslash}p{6.5cm}|}
\hline
\rowcolor{secondary!20}
\textbf{Topik} & \textbf{Solvable Problem} & \textbf{Perpetual Problem} \\
\hline
\textbf{Kecepatan Mengemudi} & Suami terburu-buru pagi karena istri lama siap-siap (situasi darurat) & Perbedaan gaya: safety vs efisiensi (terkait trust dan kontrol) \\
\hline
\textbf{Kebersihan} & Suami lupa buang sampah karena deadline kerja (situasional) & Perbedaan standar: chaos vs order (terkait nilai dan kenyamanan) \\
\hline
\textbf{Uang} & Mau beli equipment olahraga vs tabungan rumah (keputusan konkret) & Tabungan vs spending = security vs enjoyment (terkait childhood trauma) \\
\hline
\textbf{Waktu Bersama} & Suami lembur minggu ini untuk deadline (temporer) & Istri butuh quality time, suami butuh alone time (kebutuhan fundamental) \\
\hline
\textbf{Parenting} & Anak boleh main game berapa lama hari ini (keputusan spesifik) & Gaya parenting: disiplin ketat vs permisif (nilai dan keyakinan) \\
\hline
\textbf{Seks} & Belum mood malam ini karena capek (situasional) & Perbedaan libido dan definisi intimacy (kebutuhan dasar) \\
\hline
\textbf{Keluarga Besar} & Harus hadir acara keluarga minggu ini (spesifik) & Seberapa dekat dengan keluarga masing-masing (batasan dan identitas) \\
\hline
\textbf{Pekerjaan} & Butuh overtime minggu ini (temporer) & Ambisi karir vs work-life balance (prioritas hidup) \\
\hline
\end{tabular}
\end{center}

\section{Gridlock: Saat Perpetual Problem Macet dan Berbahaya}

\keyterm{Gridlock} adalah kondisi ketika perpetual problem tidak di-manage dengan baik, mengeras menjadi ``dinding'' yang memisahkan pasangan.

\subsection{Proses Terjadinya Gridlock}

\begin{center}
\begin{tikzpicture}[
    stage/.style={draw=#1, line width=2pt, rounded corners=10pt, fill=#1!10, 
                  minimum width=11cm, minimum height=1.8cm, align=left, text width=10cm}
]
    \node[stage=sagegreen] (s1) at (0,6) {
        \textbf{Stage 1: Perpetual Problem Normal}\\
        \small Pasangan mengenali perbedaan, ada konflik tapi tetap bisa dialog
    };
    
    \node[stage=accent] (s2) at (0,3.5) {
        \textbf{Stage 2: Konflik Berulang Tanpa Progress}\\
        \small Argument yang sama terjadi berulang kali, posisi semakin kaku
    };
    
    \node[stage=warmgray] (s3) at (0,1) {
        \textbf{Stage 3: Gridlock Terbentuk}\\
        \small Tidak ada lagi humor, affection, atau empathy saat membahas topik ini
    };
    
    \node[stage=primary] (s4) at (0,-1.5) {
        \textbf{Stage 4: Emotional Disengagement}\\
        \small Pasangan ``menyerah'', jarak emosional, hubungan hampa
    };
    
    \draw[->, line width=2pt] (s1) -- (s2) node[midway, right, font=\small] {Tanpa management};
    \draw[->, line width=2pt] (s2) -- (s3) node[midway, right, font=\small] {Vilification};
    \draw[->, line width=2pt] (s3) -- (s4) node[midway, right, font=\small] {Hopelessness};
\end{tikzpicture}
\end{center}

\subsection{Tanda-tanda Gridlock}

\begin{warningbox}[Tanda Gridlock yang Harus Diwaspadai]
\begin{itemize}
    \item \faExclamationTriangle\ \textbf{Konflik sama diulang berulang kali} tanpa kemajuan
    \item \faExclamationTriangle\ \textbf{Posisi semakin kaku} --- masing-masing tidak mau mengalah
    \item \faExclamationTriangle\ \textbf{Saling vilify} --- melihat pasangan sebagai ``jahat'' atau ``salah'
    \item \faExclamationTriangle\ \textbf{Tidak ada humor/affection} saat diskusi
    \item \faExclamationTriangle\ \textbf{Makin lama makin emosional} saat membahas topik
    \item \faExclamationTriangle\ \textbf{Akhirnya menghindari topik} sama sekali
    \item \faExclamationTriangle\ \textbf{Emotional disengagement} --- tidak peduli lagi
\end{itemize}

\textit{Gridlock adalah sinyal bahaya. Jika dibiarkan, hubungan akan menuju kehancuran.}
\end{warningbox}

\subsection{Dampak Gridlock pada Hubungan}

\begin{itemize}
    \item \textbf{Erosi Trust:} Kepercayaan hancur karena merasa tidak dipahami
    \item \textbf{Four Horsemen:} Criticism, Contempt, Defensiveness, Stonewalling muncul
    \item \textbf{Deadlock:} Tidak ada lagi room untuk negosiasi
    \item \textbf{Living as Roommates:} Hidup bersama tapi terpisah secara emosional
    \item \textbf{Risk of Infidelity:} Salah satu pihak mencari pemahaman di luar hubungan
\end{itemize}

\section{Dreams Within Conflict: Impian di Balik Perbedaan}

\subsection{Konsep Revolusioner Gottman}

Setiap perpetual problem menyembunyikan \keyterm{dreams} --- impian, nilai, kebutuhan, dan makna hidup yang fundamental bagi masing-masing pihak. Konflik yang tampaknya ``bodoh'' sebenarnya adalah \textbf{battle of dreams}.

\begin{insightbox}
\textit{``Di balik setiap perpetual problem, ada impian yang belum terwujud. Pasangan yang berhasil adalah yang mau mengeksplorasi dan menghormati impian pasangannya, meski impian itu bertentangan dengan impian mereka sendiri.''} --- John Gottman
\end{insightbox}

\subsection{Contoh Dreams dalam Konflik Umum}

\begin{center}
\renewcommand{\arraystretch}{1.5}
\begin{tabular}{|>{\raggedright}p{3cm}|>{\raggedright}p{4.5cm}|>{\raggedright}p{4.5cm}|>{\raggedright\arraybackslash}p{2.5cm}|}
\hline
\rowcolor{sagegreen!20}
\textbf{Topik Konflik} & \textbf{Dream Pasangan A} & \textbf{Dream Pasangan B} & \textbf{Nilai Dasar} \\
\hline
Anak & Merasa complete, legacy, kehangatan keluarga & Kebebasan, karir, spontanitas & Family vs Independence \\
\hline
Uang & Security, bebas dari anxiety masa kecil & Enjoyment, hidup sekarang & Safety vs Experience \\
\hline
Kebersihan & Order, kontrol, kenyamanan & Creativity, fleksibilitas & Structure vs Freedom \\
\hline
Waktu & Quality connection, intimacy & Personal space, recharge & Togetherness vs Alone \\
\hline
Pekerjaan & Achievement, recognition & Balance, kehadiran keluarga & Success vs Presence \\
\hline
Seks & Intimacy emosional, diterima & Physical release, stress relief & Connection vs Release \\
\hline
Rumah & Tempat istirahat yang sempurna & Tempat ekspresi kreatif & Perfection vs Expression \\
\hline
\end{tabular}
\end{center}

\subsection{Mengapa Dreams Penting}

\begin{praktikbox}[Mengapa Dreams Tidak Bisa Diabaikan]
\begin{enumerate}
    \item \textbf{Identitas:} Dreams adalah bagian dari siapa kita sebenarnya
    \item \textbf{Meaning-making:} Dreams memberi makna pada hidup
    \item \textbf{Resistence:} Mengorbankan dreams = mengorbankan diri sendiri
    \item \textbf{Conflict source:} Perbedaan dreams = sumber perpetual problem
\end{enumerate}

\textit{Solusi bukan meminta salah satu mengorbankan dreams-nya, tapi mencari \textbf{shared meaning} yang menghormati keduanya.}
\end{praktikbox}

\section{Self-Assessment: Identifikasi Jenis Masalahmu}

\begin{exercisebox}[Self-Assessment: Solvable atau Perpetual?]
\textbf{Petunjuk:} Pikirkan konflik utama yang sering terjadi dalam hubunganmu. Jawablah pertanyaan berikut dengan jujur.

\subsection*{BAGIAN A: Karakteristik Masalah}

\textbf{1. Seberapa sering konflik ini terjadi?}
\begin{itemize}
    \item[A] Jarang --- hanya dalam situasi spesifik
    \item[B] Sering --- pola yang berulang meski situasi berbeda
\end{itemize}

\textbf{2. Apakah ada solusi yang jelas?}
\begin{itemize}
    \item[A] Ya --- ada kompromi konkret yang bisa diterapkan
    \item[B] Tidak --- sepertinya tidak ada solusi yang memuaskan keduanya
\end{itemize}

\textbf{3. Seberapa intens emosi yang terlibat?}
\begin{itemize}
    \item[A] Moderate --- bisa didiskusikan tanpa terlalu emosional
    \item[B] Sangat intens --- menyentuh sesuatu yang sangat personal
\end{itemize}

\textbf{4. Apakah masalah ini terkait dengan nilai/masa kecil?}
\begin{itemize}
    \item[A] Tidak --- lebih tentang situasi sekarang
    \item[B] Ya --- berkaitan dengan pengalaman masa kecil atau nilai fundamental
\end{itemize}

\textbf{5. Bagaimana respon pasangan saat didiskusikan?}
\begin{itemize}
    \item[A] Kooperatif, mencari solusi bersama
    \item[B] Defensif, protektif, atau emosional
\end{itemize}

\subsection*{BAGIAN B: Deteksi Gridlock}

\textbf{6. Apakah kalian sering membahas topik yang sama tanpa kemajuan?}
\begin{itemize}
    \item[Ya] \textcolor{accent}{\textbf{[!]}} Tanda gridlock potensial
    \item[Tidak] Masih dalam tahap normal
\end{itemize}

\textbf{7. Apakah ada topik yang kalian hindari karena terlalu ``panas''?}
\begin{itemize}
    \item[Ya] \textcolor{accent}{\textbf{[!]}} Gridlock mungkin sudah terbentuk
    \item[Tidak] Masih bisa dialog
\end{itemize}

\textbf{8. Apakah kamu merasa pasangan ``tidak pernah mengerti''?}
\begin{itemize}
    \item[Ya] \textcolor{accent}{\textbf{[!]}} Emotional disengagement berkembang
    \item[Tidak] Masih ada harapan
\end{itemize}

\subsection*{SCORING}

\begin{center}
\renewcommand{\arraystretch}{1.4}
\begin{tabular}{|c|p{8cm}|}
\hline
\rowcolor{primary!20}
\textbf{Mayoritas Jawaban} & \textbf{Interpretasi} \\
\hline
A (BAGIAN A) & \textbf{SOLVABLE PROBLEM} --- Gunakan problem-solving skills standar \\
\hline
B (BAGIAN A) & \textbf{PERPETUAL PROBLEM} --- Fokus pada management dan dialogue \\
\hline
Ya (BAGIAN B) & \textbf{GRIDLOCK WARNING} --- Perlu intervensi segera \\
\hline
\end{tabular}
\end{center}
\end{exercisebox}

\section{Panduan Lengkap: Menyelesaikan Solvable Problems}

\subsection{5 Langkah Softened Startup}

\begin{tipbox}[Langkah 1: Softened Startup]
\begin{enumerate}
    \item \textbf{Start gentle:} ``Aku mau cerita tentang...'' bukan ``Kamu selalu...''
    \item \textbf{Complain, don't blame:} Fokus pada perilaku, bukan karakter
    \item \textbf{Describe, don't evaluate:} ``Sampah belum dibuang'' bukan ``Kamu malas''
    \item \textbf{Be polite:} Gunakan ``tolong'' dan ``terima kasih''
    \item \textbf{Give appreciations:} Mulai dengan yang positif
\end{enumerate}
\end{tipbox}

\subsection{Making and Receiving Repair Attempts}

\begin{center}
\renewcommand{\arraystretch}{1.4}
\begin{tabular}{|>{\raggedright}p{4cm}|>{\raggedright\arraybackslash}p{9cm}|}
\hline
\rowcolor{sagegreen!20}
\textbf{Tipe Repair} & \textbf{Contoh} \\
\hline
\textbf{I Feel Statement} & ``Aku merasa khawatir saat...'' \\
\hline
\textbf{Take Responsibility} & ``Aku salah, seharusnya aku...'' \\
\hline
\textbf{Apologize} & ``Maafkan aku sudah menyakitkanmu'' \\
\hline
\textbf{Agree} & ``Kamu benar tentang...'' \\
\hline
\textbf{Self-soothe} & ``Aku butuh menenangkan diri sebentar'' \\
\hline
\textbf{Affection} & Pelukan, sentuhan, humor lembut \\
\hline
\end{tabular}
\end{center}

\subsection{Step-by-Step Problem Solving}

\begin{praktikbox}[Proses 6 Langkah Problem Solving]
\textbf{Langkah 1: Define the Problem}
\begin{itemize}
    \item Masing-masing jelaskan perspektifnya (5 menit)
    \item Pasangan hanya mendengarkan, tidak interrupt
    \item Validasi: ``Jadi menurutmu...''
\end{itemize}

\textbf{Langkah 2: Brainstorm Solutions}
\begin{itemize}
    \item Buat daftar semua solusi yang mungkin
    \item Tidak ada kritik saat brainstorming
    \item Quantity over quality
\end{itemize}

\textbf{Langkah 3: Evaluate Options}
\begin{itemize}
    \item Bahas pro dan kontra masing-masing
    \item Cari solusi win-win atau kompromi
\end{itemize}

\textbf{Langkah 4: Choose a Solution}
\begin{itemize}
    \item Setujui satu solusi untuk dicoba
    \item Tentukan durasi trial period
\end{itemize}

\textbf{Langkah 5: Implement}
\begin{itemize}
    \item Jalankan solusi sesuai kesepakatan
    \item Dokumentasikan jika perlu
\end{itemize}

\textbf{Langkah 6: Follow-up}
\begin{itemize}
    \item Evaluasi setelah trial period
    \item Adjust jika perlu
\end{itemize}
\end{praktikbox}

\section{Panduan Lengkap: Mengelola Perpetual Problems}

\subsection{Paradigma Shift: From Solving to Managing}

\begin{insightbox}
\textbf{Paradoks Perpetual Problems:} Upaya untuk ``menyelesaikan'' perpetual problem justru akan memperburuknya. Strategi yang tepat adalah:
\begin{enumerate}
    \item \textbf{Accept:} Terima bahwa perbedaan ini akan selalu ada
    \item \textbf{Understand:} Pahami makna di balik posisi masing-masing
    \item \textbf{Dialogue:} Jalin dialog tentang dreams dan nilai
    \item \textbf{Compromise:} Temukan titik temu yang menghormati keduanya
    \item \textbf{Maintain:} Terus ``rawat'' dialog ini secara berkala
\end{enumerate}
\end{insightbox}

\subsection{Dreams Within Conflict Dialogue}

\begin{exercisebox}[Latihan Dreams Within Conflict]
\textbf{Tujuan:} Mengubah gridlock menjadi dialog tentang dreams

\textbf{Setup:}
\begin{itemize}
    \item Pilih perpetual problem yang belum menjadi gridlock total
    \item Pastikan kedua pihak dalam kondisi calm
    \item Siapkan 45-60 menit waktu tanpa gangguan
\end{itemize}

\textbf{Struktur Dialog:}

\textbf{Langkah 1: Speaker Shares Dreams (10 menit)}
\begin{itemize}
    \item Ceritakan tentang posisimu tanpa blame
    \item Jelaskan \textbf{mengapa} ini penting bagimu
    \item Hubungkan dengan pengalaman masa kecil/kebutuhan fundamental
    \item Contoh: ``Aku butuh banyak quality time karena di masa kecil...''
\end{itemize}

\textbf{Langkah 2: Listener Explores Dreams (10 menit)}
\begin{itemize}
    \item Tanyakan dengan rasa ingin tahu (curiosity)
    \item ``Ceritakan lebih tentang...''
    \item ``Kapan pertama kali kamu merasa ini penting?''
    \item ``Apa yang paling kamu takutkan jika kebutuhan ini tidak terpenuhi?''
\end{itemize}

\textbf{Langkah 3: Switch Roles (20 menit)}
\begin{itemize}
    \item Ulangi langkah 1-2 dengan posisi bertukar
\end{itemize}

\textbf{Langkah 4: Find Common Ground (10 menit)}
\begin{itemize}
    \item Identifikasi nilai/kebutuhan yang sama
    \item Brainstorm cara menghormati kedua dreams
    \item Buat agreement konkret
\end{itemize}
\end{exercisebox}

\subsection{Contoh Dialog Dreams Within Conflict}

\begin{center}
\begin{tikzpicture}[
    dialog/.style={draw=#1, line width=1.5pt, rounded corners=8pt, fill=#1!5,
                   minimum width=12cm, minimum height=1.2cm, align=left, text width=11cm}
]
    \node[dialog=primary] (d1) at (0,6) {
        \small\textbf{Suami (Speaker):} ``Aku tahu kita sering bertengkar soal uang. Aku suka beli gadget baru. Tapi sebenarnya, ini bukan tentang gadget...''
    };
    
    \node[dialog=warmgray] (d2) at (0,4.5) {
        \small\textbf{Istri (Listener):} ``Ceritakan lebih tentang itu...''
    };
    
    \node[dialog=primary] (d3) at (0,3) {
        \small\textbf{Suami:} ``Ayahku selalu bilang `nanti saja, kita tidak punya uang'. Aku merasa kami tidak \textbf{punya} apa-apa. Sekarang, saat aku bisa beli sesuatu, aku merasa \textbf{berhasil}. Seperti membuktikan bahwa aku bisa.''
    };
    
    \node[dialog=warmgray] (d4) at (0,1.5) {
        \small\textbf{Istri:} ``Jadi beli gadget itu tentang merasa berhasil dan membuktikan diri?''
    };
    
    \node[dialog=primary] (d5) at (0,0) {
        \small\textbf{Suami:} ``Ya. Aku takut kalau tidak bisa beli apa yang kuinginkan, aku akan kembali jadi anak kecil yang tidak punya apa-apa.''
    };
\end{tikzpicture}
\end{center}

\section{Topik Perpetual Problems Umum dan Strateginya}

\subsection{Uang: Security vs Enjoyment}

\begin{warningbox}[Konflik Uang]
\textbf{Saver (Hemat)}\ dreams: Keamanan, bebas anxiety, kontrol masa depan\\
\textbf{Spender (Habiskan)}\ dreams: Enjoyment, reward diri, hidup sekarang

\textbf{Strategi:}
\begin{itemize}
    \item Buat tiga bucket: Needs, Wants, Dreams
    \item Tetapkan persentase yang jelas untuk masing-masing
    \item ``Fun money'' --- masing-masing punya budget bebas tanpa pertanggungjawaban
    \item Review bersama setiap bulan
\end{itemize}
\end{warningbox}

\subsection{Seks: Intimacy vs Physical Release}

\begin{praktikbox}[Konflik Intimasi]
\textbf{High-desire partner}\ dreams: Connection, diterima, self-worth\\
\textbf{Low-desire partner}\ dreams: Kedekatan non-seksual, tidak merasa pressure

\textbf{Strategi:}
\begin{itemize}
    \item Scheduled intimacy --- jadwalkan tanpa terdengar ``tidak romantic''
    \item Non-sexual touch hari-hari non-intimacy
    \item Communicate needs tanpa pressure
    \item Eksplorasi apa yang membuat low-desire partner merasa dekat
\end{itemize}
\end{praktikbox}

\subsection{Parenting: Discipline vs Nurturing}

\begin{tipbox}[Konflik Parenting]
\textbf{Strict parent}\ dreams: Anak sukses, disiplin, siap dunia\\
\textbf{Lenient parent}\ dreams: Anak bahagia, kreatif, bebas

\textbf{Strategi:}
\begin{itemize}
    \item Private discussion --- jangan pernah bertengkar di depan anak
    \item United front --- kompromi dulu, baru implementasi
    \item Division of domain --- masing-masing lead di area tertentu
    \item Regular parenting meetings
\end{itemize}
\end{tipbox}

\subsection{Kebersihan: Order vs Chaos}

\begin{insightbox}[Konflik Kebersihan]
\textbf{Neat partner}\ dreams: Kontrol, kenyamanan, estetika\\
\textbf{Messy partner}\ dreams: Kreativitas, fleksibilitas, tidak obsesif

\textbf{Strategi:}
\begin{itemize}
    \item Zoning: Area bersama harus rapi, area pribadi bebas
    \item Systems, bukan blame: Buat sistem yang memudahkan
    \item Weekly reset: Waktu bersama merapikan
    \item Acceptance: Beberapa ketidaksempurnaan harus diterima
\end{itemize}
\end{insightbox}

\section{Integrasi dengan Attachment Styles}

\subsection{Bagaimana Masing-Masing Tipe Menangani Perpetual Problems}

\begin{center}
\renewcommand{\arraystretch}{1.6}
\begin{tabular}{|>{\raggedright}p{2.5cm}|>{\raggedright}p{5.5cm}|>{\raggedright\arraybackslash}p{5.5cm}|}
\hline
\rowcolor{primary!20}
\textbf{Tipe} & \textbf{Respon terhadap Perpetual Problem} & \textbf{Strategi Optimal} \\
\hline
\textbf{Anchor} & Tetap terbuka untuk dialog, tidak terlalu defensive, bisa melihat perspektif pasangan & Gunakan secure base untuk membantu pasangan merasa aman saat explore dreams \\
\hline
\textbf{Island} & Menjauh saat perpetual problem diangkat, merasa terjebak, avoid dialogue & Beri reassurance tidak akan dipaksa mengubah nilai, beri waktu proses \\
\hline
\textbf{Wave} & Emosional dan clingy saat perpetual problem muncul, fear abandonment & Validasi perasaan dulu, jangan biarkan merasa ditinggal saat dialog \\
\hline
\end{tabular}
\end{center}

\subsection{Kombinasi Tipe dan Perpetual Problems}

\begin{warningbox}[Kombinasi yang Paling Menantang]
\textbf{Island + Wave dalam Perpetual Problem:}
\begin{itemize}
    \item Saat perpetual problem muncul, Wave cling dan Island withdraw
    \item Ini memperkuat pola pursuer-distancer
    \item \textbf{Solusi:} Island harus belajar stay present, Wave harus belajar self-soothe
\end{itemize}

\textbf{Dua Island dalam Perpetual Problem:}
\begin{itemize}
    \item Keduanya withdraw, masalah tidak pernah dibahas
    \item Hubungan menjadi ``roommate'' arrangement
    \item \textbf{Solusi:} Scheduled dialogue wajib, jangan biarkan avoidance
\end{itemize}

\textbf{Dua Wave dalam Perpetual Problem:}
\begin{itemize}
    \item Konflik menjadi sangat emosional dan dramatis
    \item Testing behavior muncul
    \item \textbf{Solusi:} Grounding techniques, take turns speaking
\end{itemize}
\end{warningbox}

\section{Case Studies: Before and After}

\subsection{Case Study 1: Uang dan Gridlock}

\begin{praktikbox}[Case: David \& Sarah --- Gridlock Keuangan]
\textbf{BEFORE (Gridlock):}
\begin{itemize}
    \item David: Spender, suka beli gadget dan makan di restoran mahal
    \item Sarah: Saver, ingin tabungan rumah dan dana darurat besar
    \item Pola: Setiap kali David beli sesuatu, Sarah marah besar
    \item David mulai sembunyi-sembunyi belanja (financial infidelity)
    \item Sarah merasa tidak bisa percaya, mengontrol semua rekening
    \item Gridlock total --- tidak ada lagi dialogue
\end{itemize}

\textbf{INTERVENTI Dreams Within Conflict:}
\begin{itemize}
    \item David: ``Ayahku miskin, aku merasa tidak punya apa-apa. Beli gadget = merasa berhasil''
    \item Sarah: ``Orangtuaku bercerai karena utang. Tabungan = security aku tidak ditinggal''
\end{itemize}

\textbf{AFTER (Managed):}
\begin{itemize}
    \item Agreement: 60\% needs/tabungan, 20\% fun money masing-masing (tanpa pertanyaan), 20\% dreams
    \item David punya ``freedom account'' untuk gadget --- tidak perlu sembunyi
    \item Sarah puas dengan sistem otomatis tabungan
    \item Monthly money date untuk review dan connection
\end{itemize}
\end{praktikbox}

\subsection{Case Study 2: Waktu Bersama}

\begin{praktikbox}[Case: Mike \& Lisa --- Perpetual Time Problem]
\textbf{BEFORE:}
\begin{itemize}
    \item Lisa (Wave): Butuh quality time setiap hari, merasa diabaikan jika tidak
    \item Mike (Island): Butuh alone time setiap hari, merasa tercekik jika tidak
    \item Pola: Lisa nag $
ightarrow$ Mike withdraw $
ightarrow$ Lisa nag lebih keras $
ightarrow$ Mike menghilang
\end{itemize}

\textbf{DREAMS EXPLORED:}
\begin{itemize}
    \item Lisa: ``Di keluargaku, tidak ada yang benar-benar `melihat' aku. Quality time = aku ada dan dihargai''
    \item Mike: ``Orangtuaku selalu membutuhkan perhatianku. Alone time = aku bisa bernapas''
\end{itemize}

\textbf{AFTER:}
\begin{itemize}
    \item Structure: Mike pulang kerja $
ightarrow$ 1 jam alone time $
ightarrow$ quality time bersama
    \item Lisa belajar self-soothing saat Mike recharge
    \item Mike inisiatif quality time (bekerja sama kebutuhan Lisa)
    \item Weekend: 1 hari ``we time'', 1 hari ada alone time masing-masing
\end{itemize}
\end{praktikbox}

\section{Action Steps: Program 30 Hari}

\begin{exercisebox}[Program 30 Hari: Mastering Your Problems]

\textbf{MINGGU 1: Identification}
\begin{itemize}
    \item \textbf{Hari 1-2:} List semua konflik dalam 6 bulan terakhir
    \item \textbf{Hari 3-4:} Kategorikan: Solvable atau Perpetual? (gunakan assessment)
    \item \textbf{Hari 5-7:} Identifikasi tanda gridlock pada perpetual problems
\end{itemize}

\textbf{MINGGU 2: Solvable Problems}
\begin{itemize}
    \item \textbf{Hari 8-10:} Pilih satu solvable problem, lakukan 6-step problem solving
    \item \textbf{Hari 11-12:} Praktek softened startup dan repair attempts
    \item \textbf{Hari 13-14:} Review dan celebrate success
\end{itemize}

\textbf{MINGGU 3: Perpetual Problems --- Dreams Exploration}
\begin{itemize}
    \item \textbf{Hari 15-17:} Pilih satu perpetual problem yang belum gridlock
    \item \textbf{Hari 18-19:} Lakukan Dreams Within Conflict dialogue
    \item \textbf{Hari 20-21:} Dokumentasikan dreams masing-masing
\end{itemize}

\textbf{MINGGU 4: Integration dan Agreement}
\begin{itemize}
    \item \textbf{Hari 22-24:} Brainstorm agreement yang menghormati kedua dreams
    \item \textbf{Hari 25-26:} Implementasi trial agreement
    \item \textbf{Hari 27-28:} Check-in dan adjust
    \item \textbf{Hari 29-30:} Reflection dan planning forward
\end{itemize}
\end{exercisebox}

\section{Kesimpulan: Menerima dan Mengelola}

\begin{insightbox}
\textbf{Key Takeaway Chapter 7:}
\begin{enumerate}
    \item \textbf{69\% konflik adalah perpetual} --- ini normal, bukan tanda hubungan gagal
    \item \textbf{Solvable vs Perpetual memerlukan strategi berbeda}
    \item \textbf{Gridlock} adalah musuh sejati, bukan perpetual problem itu sendiri
    \item \textbf{Dreams within conflict} adalah kunci untuk mengelola perpetual problems
    \item \textbf{Pasangan bahagia} bukan yang tanpa konflik, tapi yang jago mengelola konflik
\end{enumerate}

\textit{``Tujuan bukan menghilangkan perbedaan, tapi belajar menari bersama perbedaan itu.''}
\end{insightbox}

\begin{tipbox}[Checklist Cepat: Apa yang Harus Diingat]
\begin{itemize}
    \item[$\square$] Cek: Apakah ini solvable atau perpetual?
    \item[$\square$] Kalau solvable: Gunakan 6-step problem solving
    \item[$\square$] Kalau perpetual: Explore dreams, jangan coba ``menyelesaikan''
    \item[$\square$] Hindari gridlock: Jaga dialogue tetap terbuka
    \item[$\square$] Integrasi dengan attachment style: Sesuaikan approach
    \item[$\square$] Review secara berkala: Perpetual problems perlu ``perawatan'' rutin
\end{itemize}
\end{tipbox}
